\subsection{A three-isogeny eliminates the unipotent residual branch}

\begin{lemma}[The normalized three-isogeny and its invariants]
\label{lem:B-RD41-isogeny}
\status{Proved; independently checked by the root and obstruction agents, including the complete proof and exact certificate.}
Over a characteristic-zero field let
\[
 E:\ y^2+a_1xy+a_3y=x^3,\qquad a_3(a_1^3-27a_3)\ne0.
\]
The quotient by the subgroup generated by $(0,0)$ is
\[
 E_1:\ y^2+a_1xy+a_3y
       =x^3-5a_1a_3x-a_3(a_1^3+7a_3).
                                                        \tag{B-RD41.1}
\]
Its invariants are
\[
 c_4(E_1)=a_1(a_1^3+216a_3),\qquad
 \Delta(E_1)=a_3(a_1^3-27a_3)^3.
                                                        \tag{B-RD41.2}
\]
For the Hesse-form coefficients $a_1=3t$, $a_3=t^3-D$, this gives
\[
          j(E_1)=\frac{27t^3(9t^3-8D)^3}{D^3(t^3-D)}.
                                                        \tag{B-RD41.3}
\]
\end{lemma}
\begin{proof}
The point $(0,0)$ has order three, and its inverse is $(0,-a_3)$,
as in \citet[Lemma~4.1, p.~8]{Fisher2015Nine}. Here the quotient
can be checked directly. On the open set $x\ne0$ set
\[
\begin{split}
 X&=x+\frac{a_1a_3}{x}+\frac{a_3^2}{x^2},\\
 Y&=y-\frac{a_3(a_1x+2a_3)y}{x^3}
        -\frac{a_3(a_1x+a_3)^2}{x^3}.
\end{split}                                             \tag{B-RD41.4}
\]
Substitution into (B-RD41.1), followed by reduction using the source
equation, gives zero identically. Also
\[
          2Y+a_1X+a_3=(2y+a_1x+a_3)\frac{dX}{dx},
\]
so the invariant differential is normalized. These are universal
polynomial identities after clearing denominators, checked by the
retained exact certificate.

The rational map extends to the smooth projective curves. The
function $X$ has pole divisor exactly
$2(O+(0,0)+(0,-a_3))$: at the two finite points $x$ is a uniformizer
and the term $a_3^2/x^2$ has nonzero leading coefficient. At $O$,
the leading term is $x$. Therefore the morphism has degree three
and sends all three displayed points to the origin. An origin-preserving
morphism of elliptic curves is a group homomorphism (apply rigidity
to the difference of the two addition maps). Its kernel is exactly
the indicated subgroup, so this is the quotient isogeny.
The standard Weierstrass invariant polynomials applied to
(B-RD41.1) give (B-RD41.2); substituting the Hesse-form coefficients
gives (B-RD41.3).
\end{proof}

\begin{proposition}[B-RD41: no level-twenty-seven failure with unipotent residual image]
\label{prop:B-RD41}
\status{Proved; independently checked by the root and obstruction agents, including the complete proof and exact certificate.}
Let $K$ be complete of mixed characteristic $(0,3)$, with algebraically
closed residue field and absolute ramification index two. No two
potentially supersingular elliptic curves with isomorphic $27$-torsion,
different prime-to-three special-inertia characters, and common
residual image $C_3$ or $C_6$ exist.

Thus any still-hypothetical height-two comparison in B.4 with
$p=3,e_K=2$ is restricted to the cyclic-four row:
\[
 \rho_3(G_K)=C_4,\qquad
 (|I_E|,|I_{E'}|)=(3,6),\qquad
 |I_{E,E'}|=18,\qquad
 v_K(j(E))=v_K(j(E'))=6.
                                                        \tag{B-RD41.5}
\]
Here $I_E,I_{E'}$ denote faithful special-inertia images, and the
pair is unordered. This proposition alone leaves least indices two
and three; Theorem~\ref{thm:B-RD42} subsequently excludes
(B-RD41.5) and proves that the least index is two.
\end{proposition}
\begin{proof}
If the common residual image is $C_6$, apply the same quadratic
twist to both curves to remove its central sign. The $27$-torsion
isomorphism persists, both $j$-invariants are unchanged, and
inequality of the two requested characters persists because both
are multiplied by the same quadratic character. We may therefore
assume the common image is the nontrivial unipotent group $C_3$.

By Proposition~\ref{prop:B-RD39}, the given comparison on nine-torsion
has the unit-parameter Hesse forms described there. After possibly
interchanging the curves, the first has parameter valuation $r=-1$
and $S=v_K(t^3-D)=-3$, and the second has $r\ge0$ and $S=0$.
The parameters $D$ and, in the reverse case, $D^{-1}$ are units.
The tame special-inertia orders are four and two.

Let $C\subset E[3]$ and $C'\subset E'[3]$ be their unique invariant
lines. These are the subgroups generated by the rational points
$(0,0)$ in the respective Hesse forms. Uniqueness follows because
a nontrivial unipotent matrix in dimension two has a unique invariant
line. The $27$-torsion isomorphism $f$ therefore carries $C$ onto $C'$.
Write $\varphi:E\to E_1=E/C$ and
$\varphi':E'\to E'_1=E'/C'$.

We spell out the level change. The inverse image of $E_1[9]$ under
$\varphi$ is
\[
 T=\{P\in E[27]:[9]P\in C\}.
\]
Indeed $[9]\varphi(P)=0$ is equivalent to $[9]P\in C$, which
also forces $[27]P=0$. Thus $E_1[9]=T/C$ as finite group schemes.
The same construction on the other side gives $T'/C'$, and $f$
induces an isomorphism
\[
                             E_1[9]\simeq E'_1[9]
                                                        \tag{B-RD41.6}
\]
over $K$. Both quotient curves are potentially supersingular.
An isogeny induces an isomorphism of rational Tate modules and of
their potential good-reduction Dieudonne modules. In particular the
special-inertia characters and their tame orders remain four and two.

Formula (B-RD41.3) gives a different pair of $j$-valuations. If
$r=-1$ and $D$ is a unit, then
$v_K(9t^3)=1$ and $v_K(8D)=0$, so
\[
                        v_K(j(E_1))=6.
\]
For the other curve, $r\ge0$ and $S=0$, while its unit parameter
is either $D$ or $D^{-1}$. Again $9t^3-8D$ (or its inverse-parameter
version) is a unit, so
\[
                       v_K(j(E'_1))=6+3r\ge6,
                                                        \tag{B-RD41.7}
\]
with $+\infty$ allowed for $t=0$.

Now classify the common residual representation in (B-RD41.6).
Since $\mu_3\subset K$, it lies in
$\operatorname{SL}_2(\mathbf F_3)$; the elementary inertia
classification in Proposition~\ref{prop:B-RD37} gives scalar image,
cyclic-four image, or cyclic-three/six image. A scalar image is
impossible by Proposition~\ref{prop:B-RD37}, since the characters
are different. Cyclic-four image is impossible by
Proposition~\ref{prop:B-RD38}, since one tame special-inertia order
is four. In the remaining case, a common quadratic twist gives
nontrivial unipotent residual image. It leaves both valuations in
(B-RD41.7) unchanged and preserves inequality of the characters.
Proposition~\ref{prop:B-RD39} would then require one $j$-valuation
to be three, contradicting (B-RD41.7). This excludes (B-RD41.6)
and hence the proposed original $27$-torsion comparison.

The reductions in Theorem~\ref{thm:B-RD34} and
Propositions~\ref{prop:B-RD35}, \ref{prop:B-RD37} identify every
remaining BT comparison with one of these elliptic residual cases.
Only the cyclic-four case survives, and
Proposition~\ref{prop:B-RD38} gives exactly (B-RD41.5).
\end{proof}

\paragraph{Exact certificate.}
The program \path{code/verify_c3_isogeny_exclusion.py} checks the
universal rational-map substitution and normalized differential,
both target invariant polynomials, the specialized quotient
$j$-formula, and all required valuation cases. It also models the
finite module passage from $27$-torsion to quotient nine-torsion,
with the actual invariant line and subgroup ranks. The retained
output records the author runtime $0.08725505400070688$ seconds.
It does not evaluate the remaining
cyclic-four antisymplectic branch.
